\documentclass[a4paper,12pt]{article}
\usepackage{listings}
\usepackage{xcolor}
\usepackage{geometry}
\usepackage{graphicx}
\usepackage{amsmath}
\usepackage{booktabs}

\geometry{margin=1in}

\definecolor{codegray}{gray}{0.9}
\lstset{
  backgroundcolor=\color{codegray},
  basicstyle=\ttfamily\small,
  frame=single,
  breaklines=true,
  keywordstyle=\color{blue},
  commentstyle=\color{green!50!black},
  numbers=left,
  numberstyle=\tiny,
  stepnumber=1,
  numbersep=8pt,
  tabsize=2,
  showspaces=false,
  showstringspaces=false,
  captionpos=b
}

\begin{document}

\begin{center}
{\LARGE \textbf{Experiment: POS-based LED Blink using VAMAN Board}} \\[6pt]
\textbf{Course: Digital Logic Design using Verilog HDL}
\end{center}

\section*{Objective}
To design and implement a Verilog module that computes a function based on its \textbf{Product of Sums (PoS)} form and uses it to control the blinking of an LED on the VAMAN FPGA board.

\section*{Apparatus Required}
\begin{itemize}
  \item VAMAN FPGA Development Board (based on Spartan-6)
  \item USB JTAG Programmer
  \item LEDs and Switches on board
  \item Xilinx ISE / Vivado Software
\end{itemize}

\section*{Logic Description}
The Boolean function is given as:
\[
f = (b + c)(a + \overline{c})
\]

The LED will blink only when the output of \( f \) is logic HIGH.  
If \( f = 0 \), the LED remains OFF.

\section*{Truth Table}
\begin{center}
\begin{tabular}{cccc|c}
\toprule
A & B & C & Expression & F \\ 
\midrule
0 & 0 & 0 & $(b+c)(a+\bar{c})$ & 0 \\
0 & 0 & 1 & $(0+1)(0+\bar{1})=1\times0$ & 0 \\
0 & 1 & 0 & $(1+0)(0+\bar{0})=1\times1$ & 1 \\
0 & 1 & 1 & $(1+1)(0+\bar{1})=1\times0$ & 0 \\
1 & 0 & 0 & $(0+0)(1+\bar{0})=0\times1$ & 0 \\
1 & 0 & 1 & $(0+1)(1+\bar{1})=1\times1$ & 1 \\
1 & 1 & 0 & $(1+0)(1+\bar{0})=1\times1$ & 1 \\
1 & 1 & 1 & $(1+1)(1+\bar{1})=1\times0$ & 0 \\
\bottomrule
\end{tabular}
\end{center}

\section*{Simplified Boolean Expression}
\[
f = (b + c)(a + \overline{c})
\]
No further simplification is required since it is already in minimized PoS form.

\section*{Verilog Code}
\begin{lstlisting}[language=Verilog, caption={Verilog code for PoS-based LED blink}]
module pos_blink(
    input wire clk,     // clock input
    input wire a,       // input a (example: connect to switch or pin)
    input wire b,       // input b (example: connect to switch or pin)
    input wire c,       // input c (example: connect to switch or pin)
    output reg led      // LED output
);

reg [23:0] cnt = 0;

// Compute minimized PoS: (b + c) & (a + ~c)
wire f;
assign f = (b | c) & (a | ~c);

// Blink logic: LED blinks only when the PoS output f is high
always @(posedge clk) begin
    if (f) begin
        cnt <= cnt + 1;
        if (cnt == 24'd5000000) begin // adjust for desired blink rate
            led <= ~led;
            cnt <= 0;
        end
    end else begin
        led <= 0;    // LED off when PoS output is low
        cnt <= 0;
    end
end

endmodule
\end{lstlisting}

\section*{Pin Configuration (VAMAN Board Example)}
\begin{center}
\begin{tabular}{lll}
\toprule
\textbf{Signal} & \textbf{FPGA Pin} & \textbf{Description} \\
\midrule
clk & P56 & On-board 50 MHz Clock \\
a & SW0 & Input Switch 0 \\
b & SW1 & Input Switch 1 \\
c & SW2 & Input Switch 2 \\
led & LED0 & Output LED 0 \\
\bottomrule
\end{tabular}
\end{center}

\section*{Conclusion}
The circuit successfully implements the given PoS expression.  
The LED blinks at a fixed rate whenever the PoS output \( f = 1 \), and remains off otherwise.

\end{document}
